\subsubsection{Biologically inspired model of synaptic \gls{wm}}

\begin{figure}[ht]
    \centering
    \includegraphics[width=\textwidth]{Figures/working-memory.png}
    \caption{Large-scale simulation of \gls{wm}. (a) Illustration of the network model. 5 different memory items (\textit{A,B,C,D} and \textit{E}) can be stored in a recurrent network of spiking neurons. Corresponding strongly connected populations within the network transiently store memory items after activation.
    (b) Network activity of the \gls{wm} model. Black dots show individual spikes of input (top) and network (bottom) neurons. Multiple phases of store and recall are shown. Insets show average firing rates (spikes per second in Hz) over recall phases. Data obtained using a current compliance of 330\,\textmu A, corresponding to average retention times of 1.5\,s, and a voltage amplitude of 0.5\,V, corresponding to a switching probability of 5\%, were used in this simulation.
    Network behavior using (c) different retention time distribution (change of compliance current) and (d) different switching probability (change of applied voltage amplitude).}
    \label{fig:tsodyks-wm}
\end{figure}

The experimental results on \gls{wm} based on volatile memristive devices can be used to develop a \gls{wm} model for simulation of a larger scale, recurrent neural network. In \cite{mongillo2008synaptic} a model of \gls{wm} in the brain is introduced that is based on \gls{stp} combined with multi-stable network dynamics. We adapted this network model to implement the core dynamics of our volatile memristive synapses (Fig.~\ref{fig:fig2} and \ref{fig:fig3}) as a model of \gls{stp}. Fig.~\ref{fig:tsodyks-wm}A illustrates our model that serves as \gls{wm} model that is able to store a set of discrete patterns. For this, the network receives input from five input populations that represent the input patterns to be stored. %The model is a variant of the synaptic \gls{wm} model that was previously introduced in \cite{mongillo2008synaptic}, adapted here to be implemented with our model of the volatile memristive synapses. 
A set of input neurons (illustrated by colored circles) are connected to the recurrent \gls{wm} network (cloud) so that the memory items could be stored and reactivated. All neurons are spiking neurons, such that outputs are given by unitary events that are emitted when an internal membrane potential variable crossed the firing threshold (see Section~\ref{sec:network}).
The network has a multi-stable dynamics in that neurons of a specific population excite each other and transiently strengthen synapses within the population through \gls{stp} when activated. Inhibitory neurons are added to facilitate the network multi-stability (see Section \textit{Materials and Methods} for details). The experimental paradigm, the network architecture, and the synapse parameters are adjusted to the characteristics of our memristive devices. Retention time parameters are fit to the measured device characteristics.

Fig.~\ref{fig:tsodyks-wm}B shows typical behavior of the model over a simulation of several seconds. Spiking activity of the network and input neurons are shown. The pattern \textit{C} is stored by strongly activating the corresponding input neurons. This triggers co-activation of the pool of \gls{wm} neurons that encode patterns \textit{C}. Through this activation, volatile memristive synapses get strengthened and cause a prominent response in \gls{wm} neurons when a recall stimulus that activates all input neurons at intermediate rates is given after a timeout of 1\,s. During this recall phase neurons that encode pattern \textit{C} show significantly increased activity. Two repeated recalls are shown and lead to reliable memory performance. In the timeout phase \gls{wm} neurons are almost perfectly silent, which is perfectly in line with the behavior of biological neurons. After the second recall phase the memory item \textit{C} is being forgotten and \textit{A} can be stored without interference by strongly activating pattern \textit{A} input neurons. In a third recall period, \textit{A} but not \textit{C} neurons get strongly activated.

The \gls{wm} model requires time constants in the order of behaviorally relevant time scales ($>$500\,ms to few seconds) to function well. Fig.~\ref{fig:tsodyks-wm}C shows the  memory recall performance of the network when different retention times are used. Recall performance is measured here as the \gls{snr} between the \gls{wm} neurons. The \gls{snr} is computed here as the mean firing rate of the population that is specific to the memory item over the mean firing rate of non-specific populations. Retention times below 1\,s lead to degraded performance because the firing activity of the neurons is significantly lower than the retention time. This result is in accordance with the experiments shown in Fig.~\ref{fig:fig3}. Also switching probabilities had to be finely tuned. Fig.~\ref{fig:tsodyks-wm}D shows the impact of switching probabilities on the recall performance. Switching probabilities of about 0.05 are found to work best for this memory store/recall protocol because it prevents an excessive activation of the volatile devices, which would prevent the correct storage of information, as also confirmed by the experimental results of Fig.~\ref{fig:fig3}. Thanks to the flexibility of our volatile memristive devices, these probabilities can be obtained with a suitable selection of the voltage amplitude and pulse time width as shown in Fig.~\ref{fig:fig2}D.

\subsubsection{Associative symbolic \gls{wm}}

\begin{figure}[ht]
    \centering
    \includegraphics[width=\textwidth]{Figures/associations.png}
    \caption{Associative symbolic \gls{wm} with volatile memristive devices. (a) Illustration of the network for associative symbolic \gls{wm}. (b) Sequence of store and recall. Store/recall input (top row) and \gls{wm} neuron output spikes (bottom row) are shown. Inserted pictograms represent the decoded objects and recall queries. Store/recall inputs had a one-to-one fixed connectivity to \gls{wm} neurons. After storing an association it can be recalled by cuing the network with an arbitrary memory element. (c) Decoding error plotted as a function of time delay between store and first recall.}
    \label{fig:symbolic-working-memory}
\end{figure}

Another important feature of \gls{wm} is the ability to transiently form associations between properties, such as color or shape, to remember representations of real-world object. Fig.~\ref{fig:symbolic-working-memory}A illustrates this associative symbolic \gls{wm} model. %Objects have here properties \emph{shape}, \emph{texture} and \emph{color} (e.g. a \textit{`smooth, red cylinder'}). 
We investigated whether the model memristive devices can be used to form transient associations between symbolic items of different categories. We used a similar store-recall paradigm as in Fig.~\ref{fig:tsodyks-wm} but here the memory item to store is given by association between features that could be dynamically bound together. These features are thought to represent the state of physical objects an associative memory networks is able to perceive through a set of sensors. Concretely we used features from 3 different stimulus categories, shape, texture and color (see Fig.~\ref{fig:symbolic-working-memory}). Objects are represented by jointly activating feature neurons in the memory network, e.g. a \textit{`smooth, red cylinder'} is represented by jointly activating the corresponding feature neurons in the network. STP model synapses as in Fig.~\ref{fig:tsodyks-wm} are used as short-term memory inside the bidirectional connections between the neurons.

The network is able to form an autoassociative memory. After an association has been formed by presenting an object to the network, the features of that object can be recalled, after a brief time delay of length $T_\text{delay}$, by querying the network with only one part of the features. Fig.~\ref{fig:tsodyks-wm}B shows one example store and recall sequence. After the object \textit{`smooth, red cylinder'} is stored, neurons representing its features can be re-activated by triggering only the \textit{`cylinder'} or \textit{`smooth'} neurons. After a new object (\textit{`smooth, blue cone'}) is stored, its representation can be retrieved by only activating the \textit{`cone'} neuron. The object representations are transiently stored for and then slowly fade away. This is analyzed in Fig.~\ref{fig:tsodyks-wm}C where we plot the decoding error as a function $T_\text{delay}$. The memory can be reliably retrieved during a time window of 600\,ms, corresponding in our devices to I$\mathrm{_{CC}}$\,=\,70\,\textmu A.
